\documentclass[a4paper]{article}

%------------------------------------------------
% الحزم
%------------------------------------------------

\usepackage{enumitem}
\usepackage[margin=2.5cm]{geometry}
\usepackage{amsmath,amsfonts,amssymb}
\usepackage[charter]{mathdesign}
\usepackage{graphicx}
\usepackage{array}
\usepackage{fancyhdr}
\usepackage{polyglossia}

%------------------------------------------------
% اللغات والخطوط
%------------------------------------------------

\setdefaultlanguage[
    calendar=gregorian,
    numerals=maghrib
]{arabic}

\setotherlanguage{english}

% إذا لم تكن الخطوط مثبتة في الجهاز، تُؤخذ من المجلد ../fonts/
\IfFontExistsTF{Amiri}{%
    \newfontfamily\arabicfont[Script=Arabic, Scale=1.15]{Amiri}
}{%
    \newfontfamily\arabicfont[
        Script=Arabic, Scale=1.15,
        Path=../fonts/, Extension=.ttf,
        UprightFont=*-Regular, BoldFont=*-Bold,
        ItalicFont=*-Italic, BoldItalicFont=*-BoldItalic
    ]{Amiri}
}

\IfFontExistsTF{AlBattar}{%
    \newfontfamily\arabicfontsf[Script=Arabic, Scale=1.2]{AlBattar}
    % خط العناوين (يعمل حتى بعد استعمال \textenglish)
    \newfontfamily\khatt[Script=Arabic, Scale=1.2]{AlBattar}
}{%
    \newfontfamily\arabicfontsf[
        Script=Arabic, Scale=1.2, Path=../fonts/,
        BoldFont=ae_AlBattar.ttf
    ]{ae_AlBattar.ttf}
    % خط العناوين (يعمل حتى بعد استعمال \textenglish)
    \newfontfamily\khatt[
        Script=Arabic, Scale=1.2, Path=../fonts/,
        BoldFont=ae_AlBattar.ttf
    ]{ae_AlBattar.ttf}
}

%------------------------------------------------
% إعداد الصفحة
%------------------------------------------------

\renewcommand{\footrulewidth}{1pt}
\renewcommand{\headrulewidth}{1pt}

\rhead{\khatt المدرسة العليا للأساتذة ببوزريعة}
\lhead{\khatt ملحقة بومرداس}

\cfoot{}
\lfoot{\khatt الأستاذ: د. بن عبد الرزاق يوسف}

\parindent 0pt
\setlength{\headsep}{15pt}
\setlength{\parskip}{3pt}

\pagestyle{fancy}

\setlist[enumerate,1]{
    label=\arabic*),
    itemsep=4pt,
    topsep=3pt
}

\setlist[itemize]{
    itemsep=2pt,
    topsep=2pt
}

\arraycolsep=3pt

%================================================
% بداية الوثيقة
%================================================

\begin{document}

%------------------------------------------------
% العنوان
%------------------------------------------------

\begin{center}

{\Large\khatt\textbf{سلسلة تمارين رقم 2}}

\vspace{0.15cm}

{\Large\khatt\textbf{فهم كيفية تعلم الذكاء الاصطناعي}}

\vspace{0.1cm}

{\khatt\textbf{البيانات ـ التدريب ـ الاختبار}}

\end{center}

\rule{\textwidth}{1pt}

\begin{center}
{\khatt\textbf{الفصل 2: فهم كيفية تعلم الذكاء الاصطناعي}}
\end{center}

\rule{\textwidth}{1pt}

%================================================
% النشاط الأول
%================================================

\section*{\khatt\underline{النشاط الأول: من البيانات إلى النموذج}}

يعتمد تعلم الآلة على البيانات لاكتشاف الأنماط والعلاقات،
ثم استعمال ما تعلمه النموذج لإنتاج توقعات أو تصنيفات جديدة.

\begin{enumerate}

%------------------------------------------------

\item اربط كل مفهوم بالتعريف المناسب:

\begin{center}

\begin{tabular}{|
>{\raggedleft\arraybackslash}m{4.2cm}|
>{\raggedleft\arraybackslash}m{10cm}|}

\hline
\textbf{المفهوم} & \textbf{التعريف} \\

\hline
بيانات التدريب
&
البيانات التي يستعملها النموذج للتعلم واكتشاف الأنماط.
\\
\hline

التدريب
&
عملية يتعلم خلالها النموذج من البيانات والعلاقات الموجودة فيها.
\\
\hline

بيانات الاختبار
&
بيانات تستعمل لتقييم النموذج بعد تدريبه.
\\
\hline

الانحياز الخوارزمي
&
إنتاج نتائج غير عادلة أو غير دقيقة بسبب مشكلات في البيانات
أو طريقة تصميم الخوارزمية.
\\
\hline

التمثيلات
&
الطريقة التي يتم بها تمثيل المعلومات داخل البيانات حتى يتمكن
النظام من معالجتها.
\\
\hline

\end{tabular}

\end{center}

%------------------------------------------------

\item لدينا مجموعة من الصور تحتوي على:

\begin{itemize}

    \item 500 صورة لهواتف ذكية.
    \item 500 صورة لحواسيب محمولة.
    \item 100 صورة لأجهزة لوحية.

\end{itemize}

يريد طالب تدريب نموذج يستطيع التمييز بين هذه الأجهزة.

\begin{enumerate}

    \item ما الذي يمكن اعتباره \textbf{بيانات التدريب}؟
    
    \item لماذا يجب أن تكون البيانات متنوعة وممثلة للحالات المختلفة؟
    
    \item ماذا يمكن أن يحدث إذا كانت أغلب الصور من نوع واحد فقط؟

\end{enumerate}

%------------------------------------------------

\item بعد تدريب النموذج، استعمل الطالب صورًا جديدة لم يستعملها
أثناء التدريب.

\begin{enumerate}

    \item ما الهدف من استعمال هذه الصور؟
    
    \item هل تعتبر هذه الصور بيانات تدريب أم بيانات اختبار؟
    
    \item ماذا يمكن أن نستنتج إذا كانت نتائج النموذج ضعيفة جدًا
    على هذه الصور؟

\end{enumerate}

%------------------------------------------------

\item أكمل السلسلة التالية:

\begin{center}

\textbf{
البيانات
$\leftarrow$
التدريب
$\leftarrow$
النموذج
$\leftarrow$
الاختبار
$\leftarrow$
النتيجة
}

\end{center}

اشرح باختصار دور كل مرحلة.

\end{enumerate}

%================================================
% النشاط الثاني
%================================================

\section*{\khatt\underline{النشاط الثاني (واجب منزلي): اكتشاف الانحياز الخوارزمي}}

\textbf{يُنجز هذا النشاط في المنزل ويُسلَّم في الحصة القادمة.}

تؤثر جودة البيانات وطريقة تمثيلها في نتائج النموذج.
وقد يؤدي نقص البيانات أو عدم تنوعها إلى نتائج غير عادلة
أو غير دقيقة.

\begin{enumerate}

\item تخيل أن جامعة درّبت نظامًا للتعرف على صور الطلبة،
لكن مجموعة التدريب تحتوي على صور لطلبة من فئة واحدة فقط.

\begin{enumerate}

    \item ما المشكلة المحتملة في بيانات التدريب؟
    
    \item ما نوع الانحياز الذي يمكن أن يظهر؟
    
    \item كيف يمكن تحسين مجموعة البيانات؟
    
\end{enumerate}

%------------------------------------------------

\item صنّف الحالات التالية إلى:

\textbf{بيانات جيدة وممثلة} أو \textbf{بيانات غير ممثلة}.

\begin{center}

\begin{tabular}{|
>{\raggedleft\arraybackslash}m{9.5cm}|
>{\centering\arraybackslash}m{5cm}|}

\hline
\textbf{الحالة} & \textbf{التصنيف} \\

\hline
مجموعة صور تحتوي على أنواع وأحجام وإضاءات مختلفة.
& \\[0.4cm]
\hline

تدريب نموذج على صور ملتقطة في ظروف إضاءة واحدة فقط.
& \\[0.4cm]
\hline

استعمال بيانات تمثل مختلف الحالات التي سيواجهها النظام.
& \\[0.4cm]
\hline

تدريب نظام على عينة صغيرة جدًا لا تمثل المستخدمين.
& \\[0.4cm]
\hline

\end{tabular}

\end{center}

%------------------------------------------------

\item اذكر طريقتين يمكن من خلالهما تقليل خطر الانحياز
الخوارزمي أثناء إعداد بيانات التدريب.

%------------------------------------------------

\item لماذا لا يكفي أن نقول:
\begin{center}
\textbf{``النموذج أعطى النتيجة، إذن النتيجة صحيحة.''}
\end{center}

علّل إجابتك.

\end{enumerate}

\newpage

\end{document}